
\subsubsection{Multi-Dimensional Trustworthiness Frameworks}

TrustLLM \citep{Sun2024} introduced an 8-dimension framework covering truthfulness, safety, fairness, robustness, privacy, machine ethics, transparency, and accountability. The framework provides per-dimension scores and model comparison capabilities across 16 mainstream LLMs. MMTrustEval proposed standardized evaluation protocols for multimodal LLMs across 5 dimensions with model cards highlighting dimension-specific performance.

These frameworks share a design limitation: dimensions are evaluated independently. TrustLLM computes aggregate scores per dimension without co-occurrence statistics showing which instances fail on multiple dimensions simultaneously. MMTrustEval model cards identify per-dimension strengths but do not quantify cross-dimensional relationships. This design reflects an implicit independence assumption---that per-dimension scores provide sufficient information for model comparison.

The present methodology differs in measurement approach. Where existing frameworks compute per-dimension accuracy, this study measures phi coefficient on instance-level binary labels across dimension pairs to quantify coupling strength when present. This reveals correlation structure that independent evaluation cannot detect by design.

\subsubsection{LLM Vulnerability Analysis}

Specialized benchmarks evaluate single dimensions. TruthfulQA measures factual accuracy on questions where models exhibit systematic falsehoods. AdvBench tests adversarial robustness via perturbations. BBQ (Bias Benchmark for QA) measures fairness across demographic groups. These provide ground truth for respective dimensions but do not capture cross-dimensional behavior.

Li and Li (2024) document triangular trade-offs between robustness and fairness---models achieving high adversarial robustness simultaneously exhibit fairness degradation in specific scenarios. This validates that coupling can exist but does not characterize breadth (how many dimension pairs couple) or model-specificity (whether different architectures show different patterns). The present methodology framework extends this observation to systematic measurement across all dimension pairs.

TrustScore \citep{Zheng2024} introduces Behavioral Consistency---a reference-free metric evaluating response alignment with intrinsic knowledge using only API outputs. This architecture-agnostic approach validates that trustworthiness properties are observable from behavior alone, supporting feasibility of coupling measurement without internal model state access.

\subsubsection{Difficulty Confound Control}

Instance difficulty confounds correlation estimates in psychometric literature. Hard instances may fail on all dimensions simultaneously, creating spurious coupling independent of shared vulnerabilities. Prior multi-dimensional LLM evaluation lacks explicit difficulty control.

This study uses partial correlation controlling for instance difficulty to isolate genuine coupling from spurious artifacts. This distinguishes shared vulnerability mechanisms (coupling persists after control) from measurement confounds (coupling disappears after control). The finding that partial phi exceeds raw phi in five of six synthetic cases (retention 86-161\%) indicates difficulty acts as suppressor variable in synthetic data by design---controlling it unmasks designed coupling strength rather than reducing artifact inflation.

\subsubsection{Positioning}

This work sits at the intersection of multi-dimensional evaluation (provides dimension coverage), behavioral detection (validates API-only measurement feasibility), and difficulty control (isolates genuine coupling from confounds). The methodology adds coupling measurement capability to existing frameworks, demonstrates measurement from behavioral outputs in synthetic validation, and validates difficulty-independence detection. The contribution is sparse coupling detection methodology demonstrated on synthetic data, not an assumption.
